\section{Related Work}
\label{sec:related_work}

This work sits at the intersection of five active research threads: (1) LLM-based unit-test generation, particularly under the mutation-testing evaluation paradigm; (2) LLM-based documentation generation from source code; (3) round-trip and consistency-based evaluation of language models; (4) LLM-as-judge protocols for semantic-equivalence assessment; and (5) multi-agent architectural specialization in generative AI systems. We survey each thread in turn and position the present study.

\subsection{LLM-based unit test generation}

The dominant paradigm for automated unit-test generation prior to 2022 was \emph{search-based software testing} (SBST): tools such as EvoSuite for Java~\citep{Fraser2011EvoSuite} and Pynguin for Python~\citep{Lukasczyk2022Pynguin} treat test-suite synthesis as an optimization problem, evolving candidate tests against a coverage- or mutation-based fitness function. These tools achieve high branch coverage on self-contained functions but suffer from the \emph{regression oracle problem} (assertions encode observed return values rather than specifications) and a \emph{per-function search-budget cost} that scales poorly to large codebases.

The advent of code-generation LLMs has shifted this landscape. \citet{Schafer2023LLMTest} demonstrated that GPT-4 can produce pytest-formatted test suites that read like hand-written tests and encode specifications drawn from docstrings or function signatures. \citet{Tufano2022LLMTest} benchmarked LLM-generated tests against SBST baselines under coverage metrics. \citet{Dakhel2024MuTAP} introduce \textsc{MuTAP}, a mutation-guided approach where surviving mutants are surfaced to the LLM as prompt-repair signals; MuTAP is the closest prior to our Path-1 (mutation kill rate) metric because both use mutation testing as the acceptance oracle for LLM-generated tests. \citet{Chen2023CodeT} and follow-up test-generation frameworks like \textsc{ChatUniTest}~\citep{Xie2024ChatUniTest} and \textsc{TestPilot}~\citep{Schafer2024TestPilot} explore LLM-based test generation with self-repair and coverage-guided feedback. \citet{Liu2024RAGSurvey} survey the retrieval-augmented-generation (RAG)~\citep{Lewis2020RAG} variants for code tasks including test generation. \citet{Dou2025LLMTestSurvey} survey the LLM-test-generation field comprehensively as of 2025. These studies broadly converge on the finding that LLM-generated tests are competitive with SBST on branch coverage while requiring no per-function search budget. Recent qualitative work by \citet{Kang2025LLMHallucinationInCode} additionally quantifies hallucination modes in LLM-generated code beyond raw accuracy metrics.

Three methodological gaps recur across this literature: (i) surface-level evaluation metrics (BLEU, ROUGE, syntactic validity) do not measure whether tests catch real defects; (ii) single-LLM studies do not surface interaction effects between LLM capability and generation method; and (iii) LLM-vs-SBST comparisons rarely appear on Python-native benchmarks with mutation-testing metrics. Building on standard mutation-testing methodology~\citep{Andrews2005MutationAnalysis, Just2014MutationAnalysis}, prior work on LLM-generated tests has produced per-model per-operator capability decompositions showing that different LLMs dominate different mutation-operator families (e.g., predicate-reasoning-tuned models on comparison and boundary operators; code-specialized MoE models on arithmetic operators). The present study takes this per-operator decomposition as its DOE oracle for per-stage strengths (Section~\ref{sec:experiment_setup_doe}).

The present study extends the mutation-testing evaluation paradigm to the \emph{cross-stage} setting: mutation kill rate is used as the closure metric for Path~1 (Eq.~\ref{eq:killrate}), and the resulting per-cell kill rates are decomposed by mutation operator to test whether single-stage per-operator capabilities transfer to the multi-stage setting (Section~\ref{sec:results_per_operator}).

\subsection{LLM-based documentation generation}

Retrieval-augmented LLM docstring generation was surveyed by \citet{Liu2024RAGSurvey} and \citet{Zhang2023RAGCode}. Prior single-stage evaluations of LLM-generated docstrings established that LLM choice matters materially for documentation quality under identical retrieval and prompting conditions --- a finding that motivates but does not answer the cross-stage question addressed here.

The present study uses docstring generation as the first stage of Path~1 ($L_{\textsf{spec}}(C) \to D'$) and as the recovered artifact in Path~3 ($L_{\textsf{spec}}(T') \to D'$). Docstring quality is evaluated via BERTScore F1 on Path~3 and, more importantly, via judge SLM equivalence rating (Section~3.5.4), because Section~\ref{sec:results_judge_metric_corr} and Section~\ref{sec:discussion_rq2} show that BERTScore alone is inadequate as a semantic-equivalence signal on the docstring--docstring closure task.

\subsection{Round-trip and consistency-based evaluation of LLMs}

Four prior studies overlap with the round-trip framing adopted here. \citet{Allamanis2024RTC} introduced \emph{Round-Trip Correctness} (RTC), an unsupervised evaluation protocol for code LLMs that composes code $\to$ natural-language $\to$ code and measures whether the reconstruction is functionally equivalent to the original. RTC is a single-LLM protocol on closed-weight models (GPT-4, Codex) and uses functional test execution as the closure metric. \citet{Zhang2025CodingTriangle} introduced the \emph{Coding Triangle} framework, evaluating triangular consistency between code, specification, and tests on competitive-programming benchmarks. The Coding Triangle study is single-LLM self-consistency, uses accuracy on hidden test suites as the closure signal, and explicitly notes ``model mixtures can enhance'' as future work. \citet{Tian2024MBCT} proposed \emph{Mutation-Based Consistency Testing} (MBCT) --- but uses mutation as the \emph{attack} on the LLM's understanding (do mutations change the LLM's output?) rather than as the closure metric (do the LLM's tests detect defects that reference-suite mutations introduce?). Most recently, \citet{Marino2026SelfConsistency} extended round-trip evaluation with a self-consistency correction that remains single-LLM.

The present study advances this line of work along five orthogonal axes:

\begin{enumerate}
  \item \textbf{Heterogeneous multi-SLM closure} as a research variable. RTC, Coding Triangle, and MBCT are single-LLM protocols; we test whether cross-LLM stage assignments differ significantly from single-LLM baselines.
  \item \textbf{Mutation kill rate as the closure metric}, adopting an SE-validated defect-detection metric~\citep{Andrews2005MutationAnalysis, Just2014MutationAnalysis} rather than surface accuracy or BLEU/ROUGE.
  \item \textbf{Open-weight coverage} with pre-registered statistical rigor: six SLMs from five distinct families, all under 30~B parameters.
  \item \textbf{Per-stage bottleneck decomposition} via leave-one-stage-out ablations (N2, N3) and mixed-effects analysis (Eq.~\ref{eq:mixedlogit}) to identify which stage owns closure success.
  \item \textbf{Human validation of the closure metric} via the 60-pair five-annotator study (three author-team + two independent, all calibrated) (Section~\ref{sec:results_human_eval}), directly quantifying the judge SLM's own agreement with human consensus ($\kappa_{\textsf{quad}} = 0.074$).
\end{enumerate}

\subsection{LLM-as-judge protocols}

\citet{Zheng2023LLMJudge} introduced the LLM-as-judge paradigm for evaluating LLM outputs, arguing that a sufficiently strong judge model can substitute for human evaluation on many tasks. Subsequent studies have documented judge-model biases including position bias, verbosity bias, and self-preference bias~\citep{Panickssery2024JudgeBias, Wang2024JudgeBias}; \citet{Wei2025BiasJudge} extend this catalogue with systematic-bias mitigation strategies. \citet{Fanous2025SycEval} introduced SycEval, a benchmark quantifying sycophancy in judge models under user pressure. The Multi-Agent Debate framework of \citet{Du2023MAD} extended judging to multi-agent ensembles, though \citet{Wynn2025DebateHarm} subsequently showed that multi-agent debate can be actively harmful when debaters share underlying RLHF biases. \citet{Chen2025LLMJudgeSurvey} provide a comprehensive 2025 survey; \citet{Li2025JudgeBench} introduce a dedicated benchmark for evaluating judge LLMs against human annotators.

The present study uses DeepSeek-R1:14B as an external single-model judge for semantic-equivalence rating on all three closure paths. DeepSeek-R1 is chosen deliberately because it belongs to a family (DeepSeek) disjoint from every pipeline SLM in the DOE, avoiding self-preference bias. Section~\ref{sec:results_judge_metric_corr} reports the Pearson correlation between judge ratings and automated metrics per path; Section~\ref{sec:discussion_rq2} stratifies the disagreement by benchmark, revealing that judge--BERTScore disagreement on Path~3 flips direction across HumanEval and MBPP --- a structural signature that supports the more general claim that BERTScore is functioning as a surface-form similarity signal rather than a semantic-equivalence signal. The judge SLM's own validity is quantified in the human-evaluation addendum (Section~\ref{sec:results_human_eval}): against the majority-vote rating from five human annotators (two of whom are independent of the pipeline runners) on the 60-pair pre-registered sample, DeepSeek-R1 attains $\kappa_{\textsf{quad}} = 0.074$ (linear-weighted $\kappa = 0.081$), placing it in Landis and Koch's ``slight'' band---a direct empirical justification for the strict-AND two-signal policy that pairs the judge with an SE-validated metric on every closure decision.

\subsection{Multi-agent architectural specialization}

The design pattern of assigning distinct roles across multiple LLM instances --- rather than encoding all responsibilities in a single monolithic prompt --- has emerged as a productive research direction in adjacent problem domains. The Constitutional AI framework of \citet{Bai2022ConstitutionalAI} argues that principle-driven AI feedback can substitute for raw human preference labels in reinforcement learning. Concurrent work by \citet{Xiong2025DelphiAgent} on DelphiAgent introduces role-differentiated agents that reach consensus through iterative rounds of independent factuality judgements. \citet{Meier2025CausalDelphi} orchestrate multiple smaller specialized models for causal discovery via structured Delphi workflows. Multi-agent orchestration frameworks such as AutoGen~\citep{Wu2024AutoGen}, MetaGPT~\citep{Hong2025MetaGPT}, and DSPy~\citep{Khattab2024DSPy} provide production tooling for such role-differentiated pipelines but do not report pre-registered DOE-style evaluations of cross-LLM stage composition on SE tasks; \citet{Li2025MultiAgentBench} survey the emerging landscape.

More recent multi-agent SE work on role separation --- for example, AgentCoder~\citep{Huang2024AgentCoder}'s programmer/test-designer/test-executor triple and Dong et al.'s~\citep{Dong2024SelfCollaboration} programmer/tester/reviewer self-collaboration --- reports that role-specialised LLM pipelines outperform monolithic single-LLM baselines on code-generation benchmarks. These studies specialise LLMs by \emph{conversational role} within a single artifact-generation task, while the present study specialises by \emph{artifact-generation stage} across the multi-artifact docstring--test--code round-trip. Neither prior line reports the pre-registered DOE-style cross-stage evaluation with a defect-detection-grounded validity metric that we contribute here.

\subsection{Positioning of this work}

The present study is, to our knowledge, the first cross-LLM $\times$ cross-stage empirical evaluation of round-trip closure in a software-engineering pipeline under a defect-detection-grounded validity metric. The five distinguishing contributions relative to the surveyed literature are:

\begin{itemize}
  \item A \textbf{pre-registered 20-cell DOE} (purposive sampling from the $6^3 = 216$ full stage-model design space; not a formal fractional factorial) covering mono (6~cells), hetero (11~cells), and null (3~cells) strata, exhaustive at the mono stratum and hypothesis-driven at the hetero stratum. The DOE is committed to source before any sweep results are observed.
  \item \textbf{Mutation kill rate as the closure metric on Path~1}, with per-mutation-operator decomposition that exposes the phi-4-in-test-stage pathology at operator resolution (Section~\ref{sec:discussion_rq3}) and empirically identifies H1 (phi-4 spec + qwen-coder test/code) as the point-estimate leader on 4 of 5 operator families (Section~\ref{sec:results_per_operator}).
  \item An \textbf{external judge SLM protocol} paired with a strict two-signal validity policy (Algorithm~2), together with post-hoc categorical decomposition (Section~\ref{sec:discussion_rq2}) that turns Path~3's aggregate judge--metric non-correlation into a specific claim about BERTScore's surface-form character.
  \item \textbf{Per-stage bottleneck attribution} via the stage-contribution figure (Fig.~\ref{fig:stage_contribution}), the capability-preservation table (Table~\ref{tab:capability_preservation}), and the leave-one-stage-out null cells N2 and~N3.
  \item An \textbf{open, tested replication package} including 67-test unit coverage (51 decision-function + 16 wiring-integration tests) of the validity decision and test-filter algorithms.
\end{itemize}
